%Text automatically generated by txt2latex
\documentclass[10pt,notitlepage]{article}
\title{txt2latex(1)}
\author{User Commands}
\date{17 May 2026}
\usepackage[margin=2cm]{geometry}
\usepackage{titlesec}
\usepackage{enumitem}
\usepackage{times}
\setlength{\parskip}{0.5\baselineskip}
\setlength{\parindent}{0pt}
\setlength{\leftskip}{3.5em}
\titlespacing*{\section}{0pt}{\baselineskip}{-0.2\baselineskip}
\setlist[itemize]{leftmargin=3.5em, labelsep=0.5em}
\setlist[enumerate]{leftmargin=3.5em, labelsep=0.5em}
\setlist[description]{labelindent=3.5em, leftmargin=7em, labelsep=1em, itemsep=0em,}
\begin{document}
\maketitle

\section*{NAME}
\textbf{txt2latex} - convert flat ASCII text to \LaTeX\  .

\section*{SYNOPSIS}
{\ttfamily
\textbf{txt2latex} [\textit{OPTION}...] \textit{FILE} 

}
\section*{DESCRIPTION}
\textbf{txt2latex} converts the input text into \LaTeX\  . 

If input file \textit{FILE} is omitted, standard input is used. 
Result is displayed on standard output. 

\textbf{txt2latex} is also able to recognize sections, paragraphs,
lists (itemize, enumerate, description), verbatim blocks as well as
inline maths and equations.

The rules for processing the text patterns are defined as following:
\begin{description}
\item[Sections] They are defined by a line in upper case starting at column 1.
                Preceding blank lines are allowed for better visualization.
\item[Paragraphs] They must be left aligned and preceded by a blank line.
                  Blank spaces at the beginning of the paragraphs are allowed
                  and there is no restriction on line alignement in a paragraph.
                  Nonetheless, identical alignment provides a better visualization
                  of the plain ASCII text. 
\item[Description list] 
                        Labels of the items are separated from the definitions
                        by at least 2 blank spaces, even before a new line, if
                        definition is too long.
\item[Itemize (bullet) list] 
                             Bullet list items are defined by the first character being "-",
                             "*" or "o" followed by a space.
\item[Enumerated list] 
                       Enumerated lists are defined by the first character being 
                       a number followed by a dot or a rounded bracket.
\item[Nested lists] 
                    Nested and mixed lists are allowed.
\item[Verbatim block] 
                      Verbatim block is used to display unmodified text such as 
                      quotes of source code.
                      They must be separated by a blank line and be indented 
                      with respect to the previuous line.
                      It will be printed using the verbatim environment.
\item[Mathematics] 
                   Inline mathematics must be enclosed with a simple \$ sign
                   and equations must be enclosed with double \$ signs.
                   For example, $E=mc^2$ is rendered as an inline math whereas 
$$ E=mc^2 $$
                   is rendered in a simple equation environment.
\item[Tables] 
              There is no support for tables.
              The workaround is put them in a verbatim block.
\end{description}

Special characters \#, \$, \%, \&, \_, \{, \}, \^\, \textbackslash\ are protected i.e. escaped.

Symbols such as $<$, $>$, $<$=, $>$= are transformed to inline math.

Special words such as \LaTeX\   and \TeX\  are turned into their equivalent 
commands in latex for generating special output.
   
\section*{OPTIONS}
\begin{description}
\item[-d date] Set date. Defaults to current date.
\item[-t mytitle] Set the title. If the title is set, txt2latex will 
                  automatically add the preambule and markups for the document
\item[-a author] Set the author.
\item[-s shift] Shift heading level by 0 (section), 1 (subsection), or 2 (subsection).
                Defaults to 0.
\item[-m, --man] Apply the conventions for man pages. See NOTES.
\item[-I txt] Italicize txt in output. Can be specified more than once.
\item[-B txt] Emphasize (bold) txt in output. Can be specified more than once.
\item[-P package] Add packages or \LaTeX\   or \TeX\  commands in the preambule.
\item[-X] Compile output with pdflatex.
\item[-v, --version] Display version.
\item[-h, --help] Display help.
\end{description}

\section*{NOTES}
The formatting rules are heavily inspired from txt2man(1). 
The special treatment for sections NAME and SYNOPSIS performed 
automatically by txt2man(1) is not done by \textbf{txt2latex}.
Nonetheless, the objective of \textbf{txt2latex} is not to parse on man-formatted
plain text but to convert any plain text to \LaTeX\  . 

Using the options -B and -I, you can easily format your text as it 
would be done by txt2man. 

\section*{EXAMPLES}
Simple conversion

\begin{verbatim}
    $ txt2latex FILE > FILE.tex
  
\end{verbatim}
Conversion with document title

\begin{verbatim}
    $ txt2latex -t "title" -a "author" -I "word" FILE > FILE.tex

\end{verbatim}
Conversion with custom packages and direct rendering with pdflatex

\begin{verbatim}
    $ txt2latex -t "title" -a "author" -I "word" -P "\\usepackage{ccfonts}" -X FILE

\end{verbatim}
Try this command to format this text itself and to produce a pdf:

\begin{verbatim}
    $ txt2latex -h 2>&1 | txt2latex -t "txt2latex(1)" -a "User commands" --man -X

\end{verbatim}
Compare it to the man page generated by txt2man(1):

\begin{verbatim}
    $ txt2latex -h 2>&1 | txt2man -s 1 -t txt2latex -v "User commands" -r 0.3.0 | man -l -Tpdf -

\end{verbatim}
\section*{SEE ALSO}
txt2man(1)
\end{document}
